\cotefile{157-4}
\scanpage{318}\mspage{42}
% archive sheet number: 317

\unsure{cause|comme}, que l'on a $\deleted{\unsure{A|B}}B(\mathrm{Fib}_{W})$, et a fortiori (comme par hypothèse $\unsure{\mathcal{F}|F}\subset\mathrm{Fib}_{W}$) $B(\mathcal{F})$.

\uline{NB} Il est clair en effet que si on a $\mathcal{F}\subset\mathcal{F}'$ alors
\[
 A(\mathcal{F}')\Longrightarrow A(\mathcal{F}) \qquad
 \text{et de \={m} pour }B,C. \tag*{(*)}
\]

Or la condition préliminaire sur $\mathcal{F}$ est bien satisfaite dans le cas \inserted{\unsure{d'une var.|d'un ens.} de l'\deleted{\unsure{une|forme}}des formes} $\deleted{\unsure{A}}\mathcal{F}=\mathcal{F}_{W_{\unsure{*|\bullet|0}}},\mathrm{Fibf}_{W_*},\mathrm{Parf}_{W_*}$. Il suffit \deleted{\unsure{des}} la voir dans le cas des sous-ensembles $\mathcal{F}_{W_*},\mathrm{Parf}_{W_*}$ de $\mathrm{Fibf}_{W_*}$, i.e que \deleted{\unsure{toute|tous}} \unsure{flèches|flèche} $f$ de $\mathrm{Cat}$ se factorise en
\[
 f=pi,
\]
avec $i\in W$, et $p$ étant, au choix, dans $\mathrm{Parf}_{W_*}$, ou dans $\mathcal{F}_{W_{\bullet}}$. Or dans le cas des $\mathrm{Parf}_{W_*}$, \margin{$\Vert$} c'est ce qui est établi dans $\overline{\underline{\mathrm{VII}}}$ (Catégories de chemins \unsure{(1)|(I)}), grâce au fonctorisme des $\underline{\mathrm{Ch}}(\unsure{\mathbb{S}|S}\unsure{.|;|,})\deleted{\unsure{g}}\longrightarrow\mathbb{S}\times\mathbb{S}$. Pour trouver une factorisation avec $f=qj$, $j\in W$, $q\in\mathcal{F}_{W_*}$, il suffit de voir que $p$ se factorise ainsi en $p=q\circ j_0$, on aura
\[
 f=\underset{\substack{\rotatebox{-90}{$\scriptstyle\in$}\\\mathcal{F}_{W_{\unsure{\bullet|0|*}}}}}{\underline{q\circ}}\,
 \underset{\substack{\rotatebox{-90}{$\scriptstyle\in$}\\W}}{\underline{(j_0i)}}\; ,
\]
\deleted{Mais si} avec \unsure{ops|obs} $f\in\mathrm{Parf}_{W_*}$. Mais si $X$ est $W$-parfait sur $S$, donc une
\scanpage{319}\mspage{43}
% archive sheet number: 318

$W$-fibré faible sur $S$, alors dans la factorisation

\diagram{\begin{tikzcd}[ampersand replacement=\&,column sep=small]
 X \arrow[rr,"i"] \arrow[dr,"f"'] \& \& \unsure{\Phi|\varphi}_{S}(\unsure{x|X}) \arrow[dl,"p"]\\
 \& S \&
\end{tikzcd}\, ,}

avec $i\in W$ et $p\in\mathrm{Catfib}_{\unsure{c|e}}$ (cf $\overline{\underline{\mathrm{VI}}}$), on sait que $p$ est aussi une \inserted{$W$-}fibration faible i.e. $p\in\unsure{\mathcal{F}|F}_{W}$, d'où la conclusion. Ainsi on trouve le

\uline{Corollaire} 4.4.2. Supposons que $\mathcal{F}$ soit l'un des six ensembles (4.4.2) (i.e. l'un des \deleted{\unsure{ens.}} des \unsure{diagrammes} \deleted{$\mathrm{Fib}$} (4.2.1), à l'exclusion de $\mathrm{Fib}_{W}$ et $\deleted{\unsure{\mathrm{Fi}}}\mathrm{Fib}_{W_{\unsure{\bullet|0|*}}}$). Alors on a
\[
 A(\mathcal{F})\Longrightarrow B(\mathcal{F})\Longrightarrow C(\mathcal{F}). \tag*{\unsure{+|*}}
\]

\uline{Lemme} 4.4.3. Soit $(M,W)$ une catégorie de modèles, et soit $\mathcal{F}\subset\mathrm{Fl}(M)$ satisfaisant les conditions suivantes

1) $\mathcal{F}$ stable par chgt de base, composition, contient les \unsure{ids|iso}

2) $M$ a un objet final, et pour \deleted{tout} $X\in\mathrm{Ob}\,M$, $X\longrightarrow e$ est dans $\mathcal{F}$

3) Toute flèche $f$ de $M$ se factorise en $f=pi$, avec $i\in W$, $p\in\mathcal{F}$ (c'est la condition du lemme \unsure{4.4.1))|4.4.1)}.

Alors les conditions $A(\mathcal{F}),B(\mathcal{F}),C(\mathcal{F})$ sur $\mathcal{F}$ sont équivalentes.

En effet, en vertu de 4.4.1, \unsure{elle|il} \deleted{\unsure{suffit que}} \inserted{reste à} \deleted{$A(\mathcal{F})\unsure{\Longrightarrow|-}$} prouver l'implication
\[
 C(\mathcal{F})\Longrightarrow A(\mathcal{F}),
\]
\scanpage{320}\mspage{44}
% archive sheet number: 319

i.e. ceci: si $\unsure{\mathcal{F}|F}$ satisfait aussi la condition

4) $\mathcal{F}\cap W\subset W^{\unsure{u|n}}$ (ou ce qui revient au \={m}, comme $\mathcal{F}$ \deleted{\unsure{satisf}} est stable par chgt de base, que $\mathcal{F}\cap W$ est stable par chgt de base),

alors on en conclut
\[
 \mathcal{F}\subset\mathrm{Fib}_{W}.
\]

Or ceci n'est autre que le lemme de K.S. Brown, sous la forme duale de celle donnée plus haut (3.2.1\unsure{;|,} \unsure{page|p.} 20).

Parmi les six ensembles de flèches envisagés dans (4.4.2), les hypothèses 1) 2) 3) sont satisfaites pour $\mathrm{Parf}_{W}$ et $\mathrm{Parf}_{W_{\unsure{0|\bullet}}}$, à l'exclusion des quatre autres. \deleted{\unsure{Une de}} \deleted{\unsure{Elles}} (les autres satisfont à \deleted{\unsure{toutes}} les conditions\unsure{;|,} à la seule exception de la stabilité par composition). On trouve donc \deleted{\unsure{provisoirement}: $\mathrm{Cat},W,W_0$}:

% slanted
\margin{\uline{NB} Pour $\mathrm{Fib}_{W_0}$, ce sont les conditions 1), 2) qui sont vérifiées, et par contre la \unsure{condition} \deleted{\unsure{2)}} 3) qui fait problème. Notons-la $F(W_0)$. Si elle est satisfaite (i.e. si \unsure{Exist|Ex} de $W_0$-fibrations), alors aussi équivalence des conditions $A(\mathrm{Fib}_{W_0})$, $B(\mathrm{Fib}_{W_0})$, $C(\mathrm{Fib}_{W_0})$.}

\uuline{Corollaire} 4.4.4. Les conditions $A(\mathrm{Parf}_{W_{\deleted{\unsure{0}}}})$, $\deleted{\unsure{A}}B(\mathrm{Parf}_{W_{\unsure{\bullet|0}}})$, $C(\mathrm{Parf}_{W})$ sont équivalentes (nous désignons cette condition sur $W$ par $D$). De même les conditions $A(\mathrm{Parf}_{W_0}),B(\mathrm{Parf}_{W_0}),C(\mathrm{Parf}_{W_0})$ \deleted{\unsure{sont équiv}} relatives
\scanpage{321}\mspage{45}
% archive sheet number: 320

(On désignera la condition obtenue sur le couple $W\supset W_0$ de loc fond. par \deleted{la condition} $D_{W_0}$ \deleted{$\unsure{S|D}$}).

Ainsi on a (comme $\mathrm{Parf}_{W_{\unsure{0|\bullet}}}\subset\mathrm{Parf}_{W}$)
\[
 D=D_W\Longrightarrow D_{W_0}, \tag*{(4.4.4)}
\]
et si $W'_0\subset W_0$, on aura \deleted{$D_{W_0}\Longrightarrow D_{W'_0}$}.
\[
 D_{W_0}\Longrightarrow D_{W'_0}. \tag*{(4.4.5)}
\]

\uline{Corollaire} 4.4.5. Désignons par $A,B,C$ respectivement les conditions $\mathrm{Loc}(7\unsure{a|\bullet}),\mathrm{Loc}(7b),\mathrm{Loc}(7\unsure{c|b})$. On a alors les implications
\[
 A\Longrightarrow B\Longrightarrow C. \tag*{(4.4.\deleted{\unsure{8}}\,\unsure{5.1|8.1})}
\]
De plus, \deleted{\unsure{si}} on désigne (pour un ens. $\unsure{\mathcal{F}|F}$ de flèches de $\mathrm{Cat}$) par $\mathrm{Comp}(\mathcal{F})$ la condition sur $\mathcal{F}$ que $\mathcal{F}$ soit stable par composition, on a les équivalences
\[
 A\Longleftrightarrow(C+\mathrm{Comp}(\mathcal{F}_{W}))
 \Longleftrightarrow(C+\mathrm{Comp}(\mathrm{Fibf}_{W})). \tag*{(4.4.\deleted{\unsure{1}}\,\unsure{5.2|5,2})}
\]

Dém. La première assertion est contenue dans le corollaire \unsure{4.4.7|4.4.2}, appliqué au cas $\mathcal{F}=\mathcal{F}_{W}$. D'autre part, si $\mathcal{F}_{W}$ satisfait à $\mathrm{Comp}(\mathcal{F}_{W})$, il satisfait aux conditions 1) \unsure{2)|?)} 3) du lemme de Brown 4.4.3, donc joint à $C$, il implique $A$. \deleted{\unsure{Pour}} Il reste, pour la première équivalence, à prouver que
\[
 A\Longrightarrow\mathrm{Comp}(\mathcal{F}_{W}).
\]
\scanpage{322}\mspage{46}
% archive sheet number: 321

Mais on a par définition,
\[
 \unsure{\mathcal{F}|F}_{W}=\mathrm{Catfib}\cap\mathrm{Fibf}_{W} \tag*{(*)}
\]
et on sait que $\mathrm{Catfib}$ est stable par composition. Donc $\mathcal{F}_{W}$ l'est aussi, pourvu que \deleted{$\mathrm{Fibf}_{W}$ le soit. Or $A$,} \inserted{\unsure{le composé|la composée} \deleted{de deux} $gf$ de} deux flèches $\in\mathcal{F}_{W}$ soit dans $\mathrm{Fibf}_{W}$. Or par hypothèse $A$, qui signifie $\mathcal{F}_{W}\subset\mathrm{Fib}_{W}$, on \unsure{aura|a} $f,g\in\mathrm{Fib}_{W}$, donc $gf\in\mathrm{Fib}_{W}$ puisque $\mathrm{Fib}_{W}$ stable par composition, a fortiori $gf\in\mathrm{Fibf}_{W}$.

D'autre part, la formule précédente $(*)$ montre que $\mathrm{Comp}(\mathrm{Fibf}_{W})\Longrightarrow\mathrm{Comp}(\deleted{\unsure{\mathrm{Cat}}}\mathcal{F}_{W})$, donc on a \inserted{gratis} l'implication $\Longleftarrow$ dans la deuxième équivalence (4.4.5.2).

L'implication en sens inverse \inserted{revient donc à} \deleted{est \unsure{moins} facile, mais elle résultera \unsure{plus ou moins|plus bas} évidente quand nous aurons prouvé plus bas
\[
 A(\mathcal{F}_{W})\Longleftrightarrow A(\mathrm{Fibf}_{W})\quad
 \text{et de \={m} pour } B,C,
\]
en sorte que $A,B,C$ peuvent s'interpréter aussi bien en termes de $\mathrm{Fibf}_{W}$, que de $\mathcal{F}_{W}$. Donc l'argument précédent \unsure{prouve|pour} aussi}
\[
 \unsure{\text{``}}A\Longrightarrow\mathrm{Comp}(\mathrm{Fibf}_{W}),
\]
\scanpage{323}\mspage{4\deleted{\unsure{6}}7}
% archive sheet number: 322

ce qui sera évident \deleted{\unsure{en}} grâce au résultat (donné plus \unsure{bas})
\[
 A=A(\unsure{\mathcal{F}|F}_{W})\Longleftrightarrow A(\mathrm{Fibf}_{W})
 \qquad\text{et de \={m} pour }B,C,
\]
qui signifie donc que \deleted{$A$}
\[
 \begin{aligned}
 \mathcal{F}_{W}\subset\mathrm{Fib}_{W}&\quad\text{implique}\quad\mathrm{Fibf}_{W}\subset\mathrm{Fib}_{W}\\
 &\hspace{4em}\text{(donc }\mathrm{Fibf}_{W}=\mathrm{Fib}_{W}),
 \end{aligned}
\]
or $\mathrm{Fib}_{W}$ étant stable par composition, il en \unsure{est} de même pour $\mathrm{Fibf}_{W}$, qed.

\uline{Corollaire} \unsure{4.4.6|4.46}. Supposons que $\mathcal{F}_{W_{\unsure{0|\bullet}}}$ soit stable par composition, ce qui est le cas notamment (\unsure{grâ à|grâce à} 4.4.5.2 appliqué à $W_0$ \unsure{ald|a.l.d.|au l de} $W$) si $W_0$ satisfait $\mathrm{Loc}(7)$. Alors les conditions $A(\mathcal{F}_{W_0}),B(\mathcal{F}_{W_0}),C(\mathcal{F}_{W_0})$ \deleted{sont équivalentes} (que nous notons aussi simplement par $A_{W_0},B_{W_0},C_{W_0}$) sont équivalentes. (\uline{NB} Si on ne fait pas l'hypothèse précédente sur $W_0$, on aura seulement $A_{W_0}\Longrightarrow B_{W_0}\Longrightarrow C_{W_0}$, en vertu de 4.4.2).

4.4.7. On a fait \unsure{le tour}, à présent de toutes les conditions \deleted{que} \inserted{\unsure{du|de} type ``fibration''} sur $W$, ou sur le couple $(W,W_0)$, qui m'ont paru utiles. Je vais les déployer ici dans un diagramme (plus ou moins) récapitulatif:
\scanpage{324}\mspage{48}
% archive sheet number: 323

\diagram{\begin{tikzcd}[ampersand replacement=\&,row sep=1ex,column sep=1em]
 \& \overset{\unsure{\mathcal{F}|F}_{W_0},\,\mathrm{Fibf}_{W_0}}{\overbrace{A_{W_0}}}
  \arrow[dd,Rightarrow] \arrow[ddddrr,Rightarrow,postaction={draw,decorate,decoration={zigzag,segment length=4pt,amplitude=2pt}},"{\deleted{\diagup}}" description] \& \& \&\\
 \overset{\mathcal{F}_{W},\,\mathrm{Fibf}_{W}}{\overbrace{A}}
  \arrow[ur,Rightarrow] \arrow[dd,Rightarrow] \& \& \& \&\\
 \& B_{W_0}\arrow[d,Rightarrow] \& \& \&\\
 B\arrow[ur,Rightarrow]\arrow[dd,Rightarrow]
 \& C_{W_0}\arrow[uuu,Rightarrow,dashed,bend right=37,"{\text{si }W_0(7)}"']
 \arrow[drr,Rightarrow,"\bullet"] \& \& \&\\
 \& \& \& \overset{\mathrm{Parf}_{W_0}}{\overbrace{D_{W_0}}}
 \arrow[r,Rightarrow,"\bullet"]
 \& \substack{F:\\(\mathrm{Cat},W,\mathrm{Fib}_{W})\\\text{une catégorie}\\\text{à fibrations},}\\
 C\arrow[uur,Rightarrow]\arrow[rr,Rightarrow,"\bullet"] \& \&
 \overset{\mathrm{Parf}_{W}}{\overbrace{D}}\arrow[ur,Rightarrow] \& \&
\end{tikzcd}}

\diagram{\begin{tikzcd}[ampersand replacement=\&,row sep=.35em,column sep=.8em]
 \& A\arrow[r,Rightarrow]\arrow[d,Rightarrow] \& A_{W_0}\arrow[d,Rightarrow] \&\\
 \mathcal{F}_{W},\,\mathrm{Fibf}_{W}\;\smash{\left\{\vphantom{\rule{0pt}{5ex}}\right.}
 \& B\arrow[r,Rightarrow]\arrow[d,Rightarrow] \& \unsure{B|D}_{W_0}\arrow[d,Rightarrow] \&
 \smash{\left.\vphantom{\rule{0pt}{5ex}}\right\}}\;\substack{\cdot\,\cdot\,\cdot\\[-.25ex]\cdot\,\cdot}\\
 \& C\arrow[r,Rightarrow]\arrow[d,Rightarrow] \& C_{W_0}\arrow[d,Rightarrow] \&\\
 \& \underset{\mathrm{Parf}_{W}}{\underbrace{D}}\arrow[r,Rightarrow]
 \& \underset{\mathrm{Parf}_{W_0}}{\underbrace{D_{W_0}}}\arrow[r,Rightarrow] \& F
\end{tikzcd}}
% slanted
\margin{version revue du diagramme}

% slanted
\margin{\deleted{3ème mouture du diagramme}}
\diagram{\begin{tikzcd}[ampersand replacement=\&,row sep=.45em,column sep=.65em,execute at end picture={\draw[overlay,thin] ([xshift=1mm]current bounding box.south west) -- ([xshift=8mm]current bounding box.north west) ([xshift=9mm]current bounding box.south west) -- ([xshift=17mm]current bounding box.north west) ([xshift=17mm]current bounding box.south west) -- ([xshift=23mm]current bounding box.north west);}]
 \deleted{A}\arrow[d,dash]\arrow[r,dash] \& \deleted{A_{W_0}}\arrow[d,dash] \& \&\\
 \deleted{B}\arrow[d,dash]\arrow[r,dash] \& \deleted{B_{W_0}}\arrow[d,dash] \& \&\\
 \deleted{C}\arrow[r,dash]\arrow[dr,dash] \& \deleted{C_{W_0}}\arrow[dr,dash] \& \&\\
 \& \deleted{D}\arrow[r,dash] \& \deleted{D_{W_0}} \& \deleted{F}\arrow[l,"{\unsure{\diagup}}" description]
\end{tikzcd}}

diagramme que je vais commenter. Toutes les implications \inserted{en traits pleins} du diagramme ont \inserted{été déjà \unsure{obtenues|établies} ou sont} tautologiques, à l'exception des trois $\overset{\bullet}{\Longrightarrow}$ sur lesquelles je vais revenir. L'implication en pointillés $C_{W_0}\dashrightarrow A_{W_0}$ est une implication conditionnelle, soumise à la condition marquée sur le diagramme comme $W_0(7)$, ce qui signifie: $W_0$ satisfait à $\mathrm{Loc}(7)$. Cette implication, i.e. l'équivalence de $A_{W_0},B_{W_0},C_{W_0}$ sous la condition $W_0(7)$, a été vue dans 4.4.6. J'ai marqué, au dessus des deux groupes $(A,B,C)$ et $(A_{W_0},B_{W_0},C_{W_0})$, ainsi qu'au dessus de $D$ et de $D'$, qui correspondent chacun à un tel groupe de conditions $A\deleted{\unsure{0}}(\mathcal{F}),B(\mathcal{F}),C(\mathcal{F})$ toutes équivalentes, avec $\mathcal{F}=\deleted{\unsure{\mathcal{F}}}\mathrm{Parf}_{W}$ ou $\mathcal{F}=\deleted{\unsure{\mathcal{F}}}\mathrm{Parf}_{W_0}$\unsure{),}, les
\scanpage{325}\mspage{49}
% archive sheet number: 324

ens de flèches $\unsure{\mathcal{F}|F}$ qui leur donnent naissance respectivement: on notera qu'il y en a deux pour les deux triplets nommés en premier, $A,B,C$ pouvant se définir via $\mathcal{F}_{W}$ ou $\mathrm{Fibf}_{W}$ indifféremment, et de \={m} $A_{W_0},B_{W_0},C_{W_0}$ via $\mathcal{F}_{W_0}$ ou $\mathrm{Fibf}_{W_0}$, comme \inserted{déjà annoncé (4.4.3) et comme} nous le verrons plus bas \inserted{(cor. 4.4.10 pages 54-55)}. Par contre, les \unsure{triplets} \deleted{\unsure{sont}} réduits $D$ et $D^{\deleted{\unsure{'}}}_{W_0}$ sont définis chacun par \deleted{$\mathrm{Parf}_{W}$} une $\mathcal{F}$ \unsure{seult|seule}, savoir par $\mathrm{Parf}_{W}$ et $\mathrm{Parf}_{W_0}$ respectivement. (En tout, cela fait intervenir six ensembles $\mathcal{F}$ distincts, à savoir les six ensembles du diagramme (4.2.1) autres que $\mathrm{Fib}_{W}$ et $\mathrm{Fib}_{W_0}$).

L'implication non tautologique $C\overset{\bullet}{\Longrightarrow}D$ est un cas particulier de l'implication $C_{W_0}\overset{\bullet}{\Longrightarrow}D_{W_0}$ \deleted{sera \unsure{prouvée} plus bas. L'implication également non tautologique $A_{W_0}\Longrightarrow D_{W_0}$ est} \inserted{laquelle} contenue dans ce qui a été annoncé \deleted{plus \unsure{bas}} \inserted{\unsure{à l'instant|d'exprimer}}, \deleted{\unsure{i.e.}} \inserted{savoir} $\deleted{\unsure{A}}\,\unsure{C|A}_{W_0}(\mathcal{F}_{W_0})\Longrightarrow \deleted{\unsure{A}}\,\unsure{C|A}_{W_0}(\mathrm{Fibf}_{W_0})$, \deleted{\unsure{i.e.}}
\deleted{\[
 \mathcal{F}_{W_0}\subset\mathrm{Fib}_{W}\ \deleted{\unsure{\Leftarrow}}\Longrightarrow\
 \mathrm{Fibf}_{W_0}\subset\mathrm{Fib}_{W}
\]}
compte tenu que $\mathrm{Parf}_{W_0}\subset\mathrm{Fibf}_{W_0}$ \inserted{cf. (\deleted{\unsure{4}}\unsure{9|*)}\unsure{.|,}\deleted{\unsure{1)|*)}} p. 42)}. \unsure{Faisant|Evidemt}
\scanpage{326}\mspage{50}
% archive sheet number: 325

$W_0=W$, on trouve donc l'implication $A\Longrightarrow D$, qui est \deleted{$A\unsure{=|\Longrightarrow}$} également contenue dans les autres implications du diagramme, comme comprise dans la chaîne d'implications
\[
 A\Longrightarrow B\Longrightarrow C\Longrightarrow D\Longrightarrow D_{W_0}
 \overset{\bullet}{\Longrightarrow}F.
\]

Il reste, dans ces commentaires généraux, à dire un mot de cette dernière implication
\[
 D_{W_0}\overset{\bullet}{\Longrightarrow}F,
\]
i.e. du fait que la plus faible de toutes les \deleted{\unsure{les|de}} conditions \deleted{\unsure{à|de} \illegible{1}} $A$ à $D_{W_0}$ envisagées ici \deleted{(à la seule exception de $B_{W_0}$ et $C_{W_0}$, qui \unsure{même})} implique déjà que $\mathrm{Cat}$ a \unsure{``}assez de $W$-fibrations, i.e est une catégorie à \unsure{cofibrations} pour $(W,\deleted{\unsure{C}}\mathrm{Fib}_{W})$. C'est clair quand on prend $D_{W_0}$ sous la forme $A(\mathrm{Parf}_{W_0})$ i.e
\[
 \unsure{\mathrm{Fi}|\mathrm{F}}\hspace{3em}\mathrm{Parf}_{W_0}\subset\unsure{\ldots W|\mathrm{Fib}_{W}}.
\]
Comme \unsure{tte|toute} flèche $f$ de $\mathrm{Cat}$ se factorise en $f=pi$, avec $\deleted{p\in}i\in W_0\subset W$ et $f\in\mathrm{Parf}_{W_0}\subset\mathrm{Fib}_{W}$, on trouve la factorisation avec $i\in W$, $f\in\mathrm{Fib}_{W}$.
\scanpage{327}\mspage{51}
% archive sheet number: 326

4.4.8. Pour terminer d'établir les implications du diagramme p. 48, il reste seulement à établir les équivalences
\[
 \left\{\begin{aligned}
 A(\deleted{\unsure{\mathrm{Fib}}}\unsure{\mathcal{F}|F}_{W_0})&\Longleftrightarrow A(\mathrm{Fibf}_{W_0})\\
 B(\mathcal{F}_{W_0})&\Longleftrightarrow B(\mathrm{Fibf}_{W_0})\\
 C(\mathcal{F}_{W_0})&\Longleftrightarrow C(\mathrm{Fibf}_{W_0}),
 \end{aligned}\right. \tag*{(4.4.8.1)}
\]
qui donneront les équivalences similaires pour $W$, en prenant $W_0=W$.

\unsure{--} La première équivalence (4.4.8.1) sera conséquence du lemme (\inserted{$\pm$} tautologique) suivant:

\uline{Lemme} 4.4.9. Soient $(M,W)$ une catégorie de modèles, $\mathcal{F},\unsure{\mathcal{G}|G}\subset\mathrm{Fl}(M)$, tels que tout $f\in\mathcal{F}$, \inserted{$f:\unsure{X|x}\longrightarrow S$} s'écrive $f=pi$, avec $i\in W^{\unsure{u|n}}_S$, $p\in\mathcal{G}$. Alors $A(\mathcal{G})\Longrightarrow A(\mathcal{F})$, i.e.
\[
 \mathcal{G}\subset\mathrm{Fib}_{W}\ \Longrightarrow\ \mathcal{F}\subset\mathrm{Fib}_{W}, \tag*{(4.4.9.1)}
\]
et de même \deleted{$B(\mathcal{G})\Longrightarrow B(\mathcal{F})$,} $C(\mathcal{G})\Longrightarrow C(\mathcal{F})$, i.e.
\[
 \mathcal{G}\cap W\subset W^u\Longrightarrow\mathcal{F}\cap W\subset W^u. \tag*{(4.4.9.2)}
\]

\uline{Dém.} A) En effet, si \deleted{\unsure{on a}} $\mathcal{G}\subset\mathrm{Fib}_{W}$, la factorisation $f=pi$ de $f$ est telle que $p\in\mathrm{Fib}_{W}$. Or un composé $pi$, avec $p\in\mathrm{Fib}_{W}$ et $i\in W_S^u$, est dans $\mathrm{Fib}_{W}$, comme on vérifie aussitôt sur le diagramme de carrés cartésiens

\diagram{\begin{tikzcd}[ampersand replacement=\&]
 X\arrow[d,dash,"{W_S^u}"',"{\unsure{-|\downarrow}}" near end] \& \unsure{\alpha'|X'}\arrow[l]\arrow[d,"{(W)}"']
 \& X''\arrow[l]\arrow[d,"{(W)}"]\\
 \overline X\arrow[d,"{\mathrm{Fib}_{W}}"'] \& \overline{X'}\arrow[l]\arrow[d,<->,"{\unsure{\uparrow|\cdot}}" near start]
 \& \overline{X''}\arrow[l,"{(W)}"]\arrow[d]\\
 S \& S'\arrow[l] \& S''\arrow[l,"W"]
\end{tikzcd}}

sur lequel on lit que $X''\unsure{\mathbin{-}|\longrightarrow}X'$ est dans $W$, qed.
\scanpage{328}\mspage{52}
% archive sheet number: 327

Ce lemme s'applique notamment au cas où $\unsure{\mathcal{F}|F}=\mathrm{Fibf}_{W_0}$, $\unsure{\mathcal{G}|G}=\mathcal{F}_{W_0}$, car $f$ se factorise en

\diagram{\begin{tikzcd}[ampersand replacement=\&]
 X\arrow[r,"i"]\arrow[d,"f"'] \& \unsure{\Phi|\varphi}_S(\unsure{x|X})\arrow[dl,"p"]\\
 S \&
\end{tikzcd}\, ,}

et on sait que $i\in W_{0\,S}^{\unsure{u|n}}\subset W_S^u$, et $p\in\mathcal{F}_{W_0}\subset\mathcal{F}_{W}$, ce qui \unsure{sont} les conditions voulues.

B) \deleted{Il faut prouver que \deleted{\unsure{nous pouvons}} dans un diagramme comm.}
\diagram{\begin{tikzcd}[ampersand replacement=\&,column sep=small,execute at end picture={\draw[overlay] (current bounding box.south west) -- (current bounding box.north east);}]
 \deleted{X}\arrow[rr,"{\deleted{u}}"]\arrow[dr,dash,"{\deleted{f}}"'] \& \& \deleted{\unsure{Y|y}}\arrow[dl,dash,"{\deleted{g}}"]\\
 \& \deleted{S} \&
\end{tikzcd}}
\deleted{les \unsure{rel.} $f,g\in\mathcal{F}$ et $u\in W$ \deleted{\unsure{sur}} impliquent \deleted{\unsure{a}} $u\in W_S^u$, sachant qu'il en va ainsi pour $\mathcal{G}$. Le \unsure{choix} de $S$.}

Démonstration similaire: Soit $f\in\mathcal{F}\cap W$, factorisons en $f=\unsure{p|g}i$, avec $i\in W_S^u$,

\diagram{\begin{tikzcd}[ampersand replacement=\&,column sep=small]
 X\arrow[dd,"{\mathcal{F}\cap W}"']\arrow[dr,dash,"{i\in W_S^u}"] \&\\
 \& \overline X\arrow[dl,dash,"g."]\\
 S \&
\end{tikzcd}}

$g\in\mathcal{G}$, on aura alors $g\in W$, donc $g\in\mathcal{G}\cap W\subset W^u$ (par hypothèse \deleted{\unsure{sur} $\mathcal{G}$} $C(\mathcal{G})$) a fortiori on a $g\in W_S^u$, donc $f\in W_S^u$ comme \unsure{composé|composée} de deux flèches de $W_S^u$, et comme le but \inserted{de $f$} est $S$ lui-même, \unsure{cel|cela} signifie $f\in W^u$, qed.

Je ne vois pas de façon \deleted{pour} prouver également
\[
 \deleted{\unsure{A}}B(\deleted{\unsure{\mathcal{F}}}\mathcal{G})\Longrightarrow B(\mathcal{F}),
\]
avec les seules hypothèses du lemme précédent sur le couple $\mathcal{F},\mathcal{G}$. Il faudra que pour tout $S$ dans $M$, et toute flèche
\scanpage{329}\mspage{53}
% archive sheet number: 328

$u:X\unsure{\mathbin{-}|\longrightarrow}Y$ dans $\deleted{\unsure{A}}M_{/S}$, on puisse insérer le \deleted{diagramme} \inserted{\unsure{triangle comm.}} obtenu \inserted{\unsure{dans $M$|ainsi}} dans un \deleted{triangle} \inserted{diagramme} \deleted{\unsure{du|de} \unsure{dit}} commutatif ci-contre

\diagram{\begin{tikzcd}[ampersand replacement=\&,column sep=small]
 X\arrow[rr,dash,"u"]\arrow[dr,"f"'] \& \&Y\arrow[dl,"g"]\\
 \& S \&
\end{tikzcd}}

avec

\diagram{\(\displaystyle(4.4.9.3)\left\{\vphantom{\rule{0pt}{17ex}}\right.\)\begin{tikzcd}[ampersand replacement=\&,column sep=small]
 X\arrow[rr,"u"]\arrow[d,"i"'] \& \&Y\arrow[d,"j"]\\
 \overline X\arrow[rr,"{\bar u}"]\arrow[dr,"p"'] \& \&\overline Y\arrow[dl,"q"]\\
 \& S \&
\end{tikzcd}}

$f=pi,g=qj,\ i,j\in W_S^{\unsure{u|n}},\ p,q\in\unsure{\mathcal{G}|G}$.

Ce sera le cas \deleted{\unsure{au}} bien \unsure{sûr|sur} si pour tt $S$, \deleted{\unsure{un $f$}} et un objet variable $X$ de $M_{/S}$, on peut trouver une factorisation \uline{fonctorielle en} $X$ de son morphisme structural $f_X$ en $p_X i_X$, avec $i_X\in W_S^u$ et $p_X\in\mathcal{G}$, i.e. si $i,p$ dérivent d'un foncteur
\[
 \unsure{\varphi|\phi}:\deleted{\unsure{M_{/S}\longrightarrow M_{/S}}}\unsure{\mathcal{F}|F}_S\longrightarrow M_{/S},
\]
% slanted
\margin{$\mathcal{F}_S$ est la souscatégorie pleine de $M_{/S}$ formée des $X$ tels que $f_X\in\mathcal{F}$.}
et d'un morphisme de foncteurs
\[
 i:\underset{\substack{\text{foncteur}\\\text{inclusion}}}{\underbrace{\deleted{\unsure{\mathrm{id}_{M_{/S}}}}\unsure{\mathrm{inc}|\mathrm{id}|j}_{\mathcal{F}_S}}}\longrightarrow\varphi.
\]

Admettant l'existence d'un \deleted{\unsure{tel}} diagramme ayant les propriétés dites, pour tte flèche $u$ dans $\mathcal{F}_S$, ou seulement tte flèche $u$ qui soit aussi dans $W_S$, nous pouvons \unsure{prouver} en outre que $u\in W_S\Longrightarrow u\in\deleted{\unsure{S}}W_S^u$, pourvu que la propriété similaire soit vraie pour $\mathcal{G}$ à la place de $\mathcal{F}$.
\scanpage{330}\mspage{54}
% archive sheet number: 329

En effet, \deleted{cette} l'hypothèse \deleted{\unsure{sur}} $u\in W$ implique $\bar u\in W$, et l'hypothèse sur $\unsure{\mathcal{G}|G}$ implique dès lors $\bar u\in W_S^{\unsure{u|n}}$. Comme $W_S^u$ est un localiseur dans $M_{/S}$, on en conclut bien $u\in W_S^u$, qed.

Ainsi on a obtenu

\deleted{Nous pouvons appliquer ce lemme et}

\uline{Corollaire} 4.4.10. Supposons, \unsure{sous les} conditions \deleted{que} du lemme 4.4.9, que tout diagramme \inserted{comm.} triangulaire

\diagram{\begin{tikzcd}[ampersand replacement=\&,column sep=small]
 X\arrow[rr,dash,"{\unsure{u|du}}"]\arrow[dr,dash,"f"'] \& \& Y\arrow[dl,dash,"g"]\\
 \& S \&
\end{tikzcd}}

s'insère dans un diagramme \inserted{du type} (4.4.9.3) comme ci-dessus, (avec les relations explicitées sur les flèches qui y figurent). Alors
\[
 \deleted{\unsure{A}}B(\mathcal{G})\Longrightarrow \deleted{\unsure{A}}B(\unsure{\mathcal{F}|F}),
\]
Donc si de plus $\mathcal{G}\subset\mathcal{F}$, \deleted{\unsure{on a l}} de sorte que $A(\mathcal{F})\Longrightarrow \deleted{\unsure{B}}A(\mathcal{G})$ \deleted{est} tautologique, et de \={m} pour $B$ et pour $C$, on trouve
\[
 \begin{aligned}
 A(\mathcal{F})&\Longleftrightarrow A(\mathcal{G})\\
 B(\mathcal{F})&\Longleftrightarrow B(\mathcal{G})\\
 C(\mathcal{F})&\Longleftrightarrow C(\mathcal{G}).
 \end{aligned}
\]

Nous pouvons appliquer ce corollaire dans le cas qui nous occupe ici, avec
\[
 \deleted{\unsure{\mathcal{F}}}\mathcal{G}\overset{\mathrm{def}}{=}\mathcal{F}_{W_0}
 \subset\mathcal{F}\overset{\mathrm{def}}{\unsure{=|\,}}\mathrm{Fibf}_{W_0},
\]
\scanpage{331}\mspage{55}
% archive sheet number: 330

vu que dans ce cas on a bien une factorisation \uline{fonctorielle} d'une $f\in\mathrm{Fibf}_{W_0}$, \inserted{$f:X\unsure{\mathbin{-}|\longrightarrow}S$} en $f=pi$, $i\in W_{0\,S}^{\unsure{u|n}}\subset W_S^u$, $p\in\unsure{\mathcal{F}|F}_{W_0}=\deleted{\unsure{\mathcal{F}}}\unsure{\mathcal{G}|G}$, \deleted{\unsure{celle qui} \unsure{savoir}} celle rappelée page 43:

\diagram{\begin{tikzcd}[ampersand replacement=\&,column sep=small]
 X\arrow[rr,"{i_X}"]\arrow[dr,"{f=f_X}"'] \& \&
 \deleted{\unsure{Z|X}}\overline X=\unsure{\Phi|\varphi}_S(\unsure{X|x})\arrow[dl,"{p_X}"]\\
 \& S \&
\end{tikzcd}}

Cela prouve donc les équivalences annoncées (4.4.8.1) et achève d'établir la validité des implications du diagramme de la page 48.

\deleted{\unsure{4}} \deleted{\unsure{Nouvelles}} \inserted{Ces implications}, à la seule exception de l'implication en pointillés $C_{W_0}\mathrel{\begin{tikzpicture}[baseline=-.5ex]\draw[double,dashed,-{Implies}] (0,0)--(.55,0);\end{tikzpicture}}A_{W_0}$, \deleted{\unsure{ont}} (qui suppose qu'on ait $W_0(7)$), ne font appel à aucune hypothèse particulière sur les localisateurs fondamentaux $W$ et $W_0$ (à part l'inclusion $W\supset W_0$).

Mais nous pouvons, en utilisant le lemme 4.4.10, donner des implications conditionnelles remarquables, sous la seule condition que $W$, ou $W_0$, ou $W$ et $W_0$, satisfassent à l'un des \uline{axiomes}
\scanpage{332}\mspage{56}
% archive sheet number: 331

\uline{des limites} $\mathrm{Loc}(8)$, et \={m} seulement sous la forme faible de stabilité \deleted{\unsure{de $W$}} \unsure{du localiseur|des localiseurs} par \uline{lim} de \uline{suites}. Donnons-les d'abord sous forme technique, avant de déployer les conséquences en un théorème-synthèse:

\uline{Corollaire} 4.4.11. Supposons que $W_0$ satisfasse \unsure{à|!} l'axiome $\mathrm{Loc}(8)$, ou seulement $\mathrm{Loc}(8\,\mathrm{bis})$ (stabilité par \uline{lim} de suites). Alors $D_{W_0}\unsure{\Longleftrightarrow|\Longrightarrow}A_{W_0}$ \inserted{(donc les conditions $A_{W_0},B_{W_0},C_{W_0},D_{W_0}$ sont équivalentes)}. En particulier (faisant $W_0=W$), si $W$ satisfait à $\mathrm{Loc}(8\,\mathrm{bis})$, alors on a $D_W\Longleftrightarrow A_W$. (donc les conditions $A_W,B_W,C_W,D_W$ sont équivalentes).

En effet, l'implication à prouver s'écrit aussi
\[
 A(\mathrm{Parf}_{W_0})\Longleftrightarrow A(\mathrm{Fibf}_{W_0}).
\]
\deleted{(l'implication $\Longleftarrow$ étant tautologique). \unsure{Cette}}

L'implication\unsure{.} \inserted{\unsure{or|?}} non tautologique $\Longrightarrow$ est justiciable du lemme 4.4.9. Il faut \deleted{\unsure{prouver}} \inserted{vérifier} que, \unsure{moyennant} $\mathrm{Loc}(8\,\mathrm{bis})$ pour $W_0$, toute $f\in\unsure{\mathrm{Fib}|\mathrm{Fibf}}_{W_0}$, \inserted{$f:X\unsure{\longrightarrow|\mathbin{-}}S$,} se factorise en $f=pi_{\unsure{*|x}}$, avec $i\in W_S^{\unsure{u|n}}$, $p\in\mathrm{Parf}_{W_0}$. Or on sait (VI, p. \unsure{193|193$_1$} ff)

% slanted
\margin{et c'est là \unsure{une trouvaille} tout à fait inattendue, \unsure{presque|à mes yeux} comme le lemme K. S. Brown\unsure{\ldots}}
